% Slide 4 -- Part B2: NVIDIA MPS (Multi-Process Service).
\begin{frame}[t]{NVIDIA MPS: Thực thi song song nhiều tiến trình CUDA}
\begin{canvas}
\node[chip,anchor=north west,inner ysep=2.5pt,font=\fontsize{8.2}{10}\selectfont\bfseries] at (0,6.28){KỸ THUẬT 2\quad /\quad NVIDIA MULTI-PROCESS SERVICE (MPS)};

% Top diagram: Clients -> MPS Server -> Concurrent GPU Execution
\node[scopebox,draw=Teal!45,fill=Teal!2,minimum width=14.55cm,minimum height=2.04cm,anchor=north west] (mpsframe) at (0,5.78){};

\node[smallbox,fill=BK!8,draw=BK!45,text width=1.85cm,align=center,minimum height=.44cm,inner sep=3pt,font=\fontsize{8.4}{10}\selectfont\bfseries] (ca) at (1.25,5.38){CUDA Client A};
\node[smallbox,fill=Teal!10,draw=Teal!50,text width=1.85cm,align=center,minimum height=.44cm,inner sep=3pt,font=\fontsize{8.4}{10}\selectfont\bfseries] (cb) at (1.25,4.76){CUDA Client B};
\node[smallbox,fill=Amber!10,draw=Amber!50,text width=1.85cm,align=center,minimum height=.44cm,inner sep=3pt,font=\fontsize{8.4}{10}\selectfont\bfseries] (cc) at (1.25,4.14){CUDA Client C};

\node[smallbox,fill=white,draw=Teal!50,text width=3.1cm,minimum height=1.56cm,align=center,inner sep=6pt] (mpssrv) at (5.35,4.76){%
 \textcolor{Teal}{\textbf{MPS Server}}\\[.08cm]
 {\fontsize{8}{9.8}\selectfont Triển khai CUDA API\\tương thích nhị phân\\Gộp luồng lập lịch}};

\node[smallbox,fill=white,draw=BK!35,text width=5.25cm,minimum height=1.66cm,anchor=north west,inner sep=6pt] (mpsgpu) at (8.55,5.59){};
\node[anchor=north west,text=BK,font=\fontsize{8.8}{10.5}\selectfont\bfseries] at (8.74,5.47){1 GPU vật lý --- Chạy song song};
\fill[BK!75,rounded corners=2pt] (8.76,4.90) rectangle (13.98,5.16);
\node[anchor=center,text=white,font=\fontsize{7.8}{9}\selectfont\bfseries] at (11.37,5.03){Kernels Client A (SMs)};
\fill[Teal!75,rounded corners=2pt] (8.76,4.54) rectangle (13.98,4.80);
\node[anchor=center,text=white,font=\fontsize{7.8}{9}\selectfont\bfseries] at (11.37,4.67){Kernels Client B (SMs)};
\fill[Amber!80,rounded corners=2pt] (8.76,4.18) rectangle (13.98,4.44);
\node[anchor=center,text=white,font=\fontsize{7.8}{9}\selectfont\bfseries] at (11.37,4.31){Kernels Client C (SMs)};

\draw[arr] (ca.east)--(mpssrv.west |- ca.east);
\draw[arr,draw=Teal] (cb.east)--(mpssrv.west);
\draw[arr,draw=Amber] (cc.east)--(mpssrv.west |- cc.east);
\draw[arr,draw=Teal] (mpssrv.east)--(mpsgpu.west |- mpssrv.east);

% Middle 2-column structured cards
\node[smallbox,fill=Teal!4,draw=Teal!35,text width=6.5cm,minimum height=2.40cm,anchor=north west] (mpspro) at (0,3.48){};
\node[anchor=north west,text=Teal,font=\fontsize{10.2}{12}\selectfont\bfseries] at (.22,3.30){Nguyên lý \& Điểm mạnh};
\node[anchor=north west,text=Ink,text width=6.25cm,font=\fontsize{8.7}{11.4}\selectfont] at (.22,2.78){%
 $\bullet$\quad Các CUDA kernel chạy \textbf{đồng thời} để lấp đầy SM của GPU.\\
 $\bullet$\quad Giảm overhead context switching và bộ nhớ ngữ cảnh.\\
 $\bullet$\quad GKE giới hạn mềm \textbf{\% active threads} \& \textbf{pinned memory}.};

\node[smallbox,fill=Red!5,draw=Red!35,text width=6.5cm,minimum height=2.40cm,anchor=north west] (mpscon) at (7.55,3.48){};
\node[anchor=north west,text=Red,font=\fontsize{10.2}{12}\selectfont\bfseries] at (7.77,3.30){Hạn chế \& Rủi ro cô lập};
\node[anchor=north west,text=Ink,text width=6.25cm,font=\fontsize{8.7}{11.4}\selectfont] at (7.77,2.78){%
 $\bullet$\quad Không cô lập phần cứng cho memory bandwidth/encoder.\\
 $\bullet$\quad \textbf{Lỗi nghiêm trọng ở 1 client} lan sang mọi client trên GPU.\\
 $\bullet$\quad GKE yêu cầu \texttt{hostIPC: true} (cần chú ý bảo mật).};

\auditnode{cards}{mpssrv}
\auditnode{cards}{mpsgpu}
\auditnode{cards}{mpspro}
\auditnode{cards}{mpscon}
\keynote{Phù hợp}{Batch nhỏ/vừa hoặc các tiến trình hợp tác (cooperative MPI/CUDA) cần tối đa hóa throughput.}
\end{canvas}
\end{frame}
